\documentclass[12pt]{article}
\usepackage[margin=0.6in,top=1.15in,bottom=0.55in]{geometry}
\usepackage{lmodern}
\usepackage{microtype}
\usepackage{titlesec}
\usepackage{enumitem}
\usepackage{needspace}
\usepackage[hidelinks]{hyperref}
\usepackage{../../latex/grader-worksheet}

\setlength{\parindent}{0pt}
\setlength{\parskip}{0.25em}
\titleformat{\section}{\large\bfseries\scshape}{}{0pt}{}

\GraderSetup{assignment-id=U1L6BLE,template-revision=1}

\begin{document}

\begin{center}
{\LARGE\bfseries Week 6B Lit Example Worksheet}\par\vspace{0.05em}
{\large\scshape Compound Pressure in Caesar and Macbeth}\par\vspace{0.12em}
\rule{0.78\textwidth}{0.5pt}
\end{center}

\textbf{Purpose:} Apply Week 6B compound tools to Julius Caesar and Macbeth. Separate either/or and if/then form from fact and persuasion.

\textbf{Directions:} Use the Lit Example Reader. Do not treat emotional force as proof.

\section*{Part 1: Brutus's Either/Or}

\textbf{Question 1.1}\\[-0.15em]
State Brutus's either/or choice in plain English. Identify both branches.

\GraderAnswerBox{Part 1 Q1}{2.25in}

\textbf{Question 1.2}\\[-0.15em]
Name at least one live option the either/or leaves out. Why does that omitted option matter for a careful listener?

\GraderAnswerBox{Part 1 Q2}{2.35in}

\newpage

\section*{Part 2: Loaded Terms and Fact Pressure}

\textbf{Question 2.1}\\[-0.15em]
Which word or pair of words in Brutus's speech most needs to be fixed before the either/or can fairly force a choice? Explain.

\GraderAnswerBox{Part 2 Q1}{2.35in}

\textbf{Question 2.2}\\[-0.15em]
Even if the either/or form were cleaned up, what factual claim in Brutus's speech would still need proof before the killing is justified to a careful Roman?

\GraderAnswerBox{Part 2 Q2}{2.35in}

\newpage

\section*{Part 3: Macbeth Conditionals}

\textbf{Question 3.1}\\[-0.15em]
Reconstruct one if/then pressure from Passage 2. Label the antecedent and the consequent.

\GraderAnswerBox{Part 3 Q1}{2.45in}

\textbf{Question 3.2}\\[-0.15em]
A character treats the thane of Cawdor title as support for the greater promise of kingship. Explain why affirming a later success does not by itself prove that the whole prophecy chain is established.

\GraderAnswerBox{Part 3 Q2}{2.55in}

\newpage

\section*{Part 4: Final Diagnosis}

\textbf{Question 4.1}\\[-0.15em]
Compare the two passages: name the main compound form in each (either/or or if/then) and state one thing that remains unproved in each case.

\GraderAnswerBox{Part 4 Q1}{2.55in}

\textbf{Question 4.2}\\[-0.15em]
In one short paragraph, explain why a vivid compound claim can feel necessary without being sound. Use one detail from Caesar or Macbeth.

\GraderAnswerBox{Part 4 Q2}{2.45in}

\end{document}
